\documentclass[../Odyssey_Project.tex]{subfiles}

\begin{document}
\setcounter{section}{6}
\section{Sylow's Theorem}

\subsubsection*{Setup and notation}
\begin{itemize}
    \item Throughout the section, $G$ is a finite group and $p$ a prime divisor of $|G|$.
    \item $|G|_p$ denotes the highest power of $p$ dividing $|G|$, so $|G|_p = p^n$ with $p^n \mid |G|$ but $p^{n+1} \nmid |G|$.
\end{itemize}

\subsubsection*{$p$-elements, $p$-groups, $p$-subgroups}
\begin{itemize}
    \item $g \in G$ is a \emph{$p$-element} if its order is a power of $p$.
    \item $G$ is a \emph{$p$-group} if $|G|$ is a power of $p$; $H \leq G$ is a \emph{$p$-subgroup} if $|H|$ is a power of $p$.
    \item Every element of a $p$-group is a $p$-element. (For infinite groups, $p$-group means every element is a $p$-element.)
    \item $H \leq G$ is a \emph{Sylow $p$-subgroup} of $G$ if $H$ is a $p$-subgroup of order $|G|_p$, the maximal possible order of a $p$-subgroup of $G$.
    \item A subgroup of $G$ is a \emph{Sylow subgroup} if it is a Sylow $p$-subgroup for some prime divisor $p$ of $|G|$.
\end{itemize}

\subsubsection*{Motivating example: $\mathrm{GL}(n, p)$}
\begin{itemize}
    \item For $G = \mathrm{GL}(n, p)$, Proposition~\ref{prop:4.2} gives $|G| = p^{\frac{n(n-1)}{2}}(p^n - 1)\cdots(p - 1)$, so $|G|_p = p^{\frac{n(n-1)}{2}}$.
    \item The subgroup $U$ of upper uni-triangular matrices satisfies $|U| = p^{n-1}p^{n-2}\cdots p^2 p = p^{\frac{n(n-1)}{2}} = |G|_p$, hence $U$ is a Sylow $p$-subgroup.
    \item Every $p$-element of $\mathrm{GL}(n,p)$ is unipotent (its characteristic polynomial is $(X-1)^n$), and by \nameref{thm:kolchin} any $p$-subgroup is conjugate with a subgroup of $U$.
    \item In particular: $\mathrm{GL}(n,p)$ has a Sylow $p$-subgroup, all are conjugate, and every $p$-subgroup is contained in one. \nameref{thm:sylow} extends this to arbitrary finite $G$.
\end{itemize}

\begin{namedtheorem}[Sylow's Theorem]
\label{thm:sylow}
(i) $G$ has at least one Sylow $p$-subgroup.

(ii) All the Sylow $p$-subgroups of $G$ are conjugate.

(iii) Any $p$-subgroup of $G$ is contained in a Sylow $p$-subgroup.

(iv) The number of Sylow $p$-subgroups of $G$ is congruent to $1$ modulo $p$.
\end{namedtheorem}
\begin{proof}
Let $|G| = p^n m$, where $p^n = |G|_p$ and $p\nmid m$. Let $X$ be the set of all subsets of $G$ having $p^n$ elements. Then $G$ acts on $X$ by left multiplication.\\
We first examine the orbits of this action whose sizes are not divisible by $p$. Suppose that $\mathcal{O}$ is such an orbit. Choose some $A \in \mathcal{O}$. Since $A$ is non-empty, there is some $a \in A$, and then $a^{-1}A$ lies in the same orbit and contains $1$. Replacing $A$ by $a^{-1}A$, we may assume that $1 \in A$. Let $P$ be the stabilizer of $A$. If $x \in P$, then $xA = A$, so $x = x1 \in A$. Hence $P \subseteq A$, and therefore
\[
    |P| \leq |A| = p^n.
\]
By Corollary~\ref{cor:3.5}, we have
\[
    |\mathcal{O}| = |G : P|,
\]
and hence $|G| = |P|\,|\mathcal{O}|$. Since $p\nmid |\mathcal{O}|$, the full $p$-part $p^n$ of $|G|$ must divide $|P|$. Together with $|P|\leq p^n$, this gives $|P|=p^n$. Thus $P$ is a Sylow $p$-subgroup of $G$. Since $P \subseteq A$ and $|P|=|A|$, we have $A=P$, so $\mathcal{O}$ is the coset space $G/P$.\\
Conversely, if $P$ is a Sylow $p$-subgroup of $G$, then the set of left cosets $G/P$ is contained in $X$, and it is exactly the orbit of the subset $P$ under left multiplication. This orbit has size
\[
    |G : P| = m,
\]
which is not divisible by $p$. Therefore the Sylow $p$-subgroups of $G$ are in bijective correspondence with the orbits in $X$ whose sizes are not divisible by $p$.\\
Let $X'$ be the union of all orbits in $X$ whose sizes are not divisible by $p$, and let $r$ be the number of Sylow $p$-subgroups of $G$. Every orbit in $X-X'$ has size divisible by $p$, so
\[
    |X| \equiv |X'| \pmod{p}.
\]
Each orbit contained in $X'$ has size $m$, and there are $r$ such orbits. Hence
\[
    |X'| = rm,
\]
so
\[
    rm \equiv |X| = \binom{p^n m}{p^n} \pmod{p}.
\]
This congruence depends only on the order of the group. Now compare with the cyclic group of order $p^n m$. By Theorem~\ref{thm:1.4}, this cyclic group has exactly one subgroup of order $p^n$, so the same congruence gives
\[
    m \equiv \binom{p^n m}{p^n} \pmod{p}.
\]
Since $p\nmid m$, we may cancel $m$ modulo $p$, obtaining
\[
    r \equiv 1 \pmod{p}.
\]
This proves (iv), and in particular $r>0$, so $G$ has at least one Sylow $p$-subgroup. This proves (i).\\
It remains to prove conjugacy and containment. Let $P$ be a Sylow $p$-subgroup of $G$, and let $Q$ be an arbitrary $p$-subgroup of $G$. Let $Y$ be the set of conjugates of $P$ in $G$. The group $Q$ acts on $Y$ by conjugation. Each orbit of this action has size equal to the index of a subgroup of the $p$-group $Q$, and hence has size a power of $p$.\\
By Proposition~\ref{prop:3.14},
\[
    |Y| = |G : N_G(P)|.
\]
Since $P \leq N_G(P)$, the index $|G : N_G(P)|$ divides
\[
    |G : P| = m.
\]
Thus $p\nmid |Y|$. Therefore at least one orbit of the action of $Q$ on $Y$ has size $1$. Let this orbit be $\{P_1\}$, where $P_1$ is a conjugate of $P$. Then $xP_1x^{-1}=P_1$ for every $x\in Q$, so $QP_1=P_1Q$. By Proposition~\ref{prop:1.3}, $QP_1\leq G$. Moreover, by Proposition~\ref{prop:3.12},
\[
    |QP_1|=\frac{|Q||P_1|}{|Q\cap P_1|},
\]
so $QP_1$ is a $p$-subgroup of $G$. Since $P_1$ is already a Sylow $p$-subgroup and $P_1\leq QP_1$, we must have $QP_1=P_1$. Hence $Q\leq P_1$. This proves that every $p$-subgroup of $G$ is contained in a Sylow $p$-subgroup, proving (iii).\\
Finally, if $Q$ itself is a Sylow $p$-subgroup of $G$, then $Q\leq P_1$ and $|Q|=|P_1|$, so $Q=P_1$. Since $P_1$ is a conjugate of $P$, every Sylow $p$-subgroup of $G$ is conjugate to $P$. This proves (ii).
\end{proof}

\begin{corollary}
\label{cor:7.1}
The number of Sylow $p$-subgroups of $G$ divides $|G|/|G|_p$.
\end{corollary}
\begin{proof}
Let $r$ be the number of Sylow $p$-subgroups of $G$, and let $P$ be one of them. By part (ii) of \nameref{thm:sylow}, every Sylow $p$-subgroup of $G$ is conjugate to $P$, so $r$ is exactly the number of conjugates of $P$ in $G$. By Proposition~\ref{prop:3.14}, this number is
\[
    r = |G : N_G(P)|.
\]
Since $P \leq N_G(P) \leq G$, we have
\[
    |G : P| = |G : N_G(P)|\,|N_G(P) : P|.
\]
Therefore $r=|G:N_G(P)|$ divides
\[
    |G:P|=\frac{|G|}{|P|}=\frac{|G|}{|G|_p},
\]
as required.
\end{proof}

\begin{namedtheorem}[Cauchy's Theorem]
\label{thm:cauchy}
$G$ has an element of order $p$.
\end{namedtheorem}
\begin{proof}
By part (i) of \nameref{thm:sylow}, $G$ has a Sylow $p$-subgroup $P$. Since $p$ divides $|G|$, this subgroup is non-trivial. Choose some $1\neq x\in P$. Then $o(x)$ is a power of $p$, say $o(x)=p^k$ with $k\geq 1$. Therefore
\[
    o\left(x^{p^{k-1}}\right)=p.
\]
Thus $G$ has an element of order $p$.
\end{proof}

\noindent (Other proofs of \nameref{thm:sylow} often assume that \nameref{thm:cauchy} is already known, as was the case historically.)

\subsubsection*{Sylow subgroups of normal subgroups and quotients}
\begin{itemize}
    \item Sylow $p$-subgroups behave well with respect to a normal subgroup $N \trianglelefteq G$: both $PN/N$ and $P \cap N$ are Sylow $p$-subgroups (of $G/N$ and $N$ respectively).
    \item For a non-normal subgroup $H \leq G$, $P \cap H$ need not be Sylow in $H$, but some conjugate $gPg^{-1} \cap H$ is.
\end{itemize}

\begin{proposition}
\label{prop:7.2}
Let $N \trianglelefteq G$, and let $P$ be a Sylow $p$-subgroup of $G$. Then $PN/N$ is a Sylow $p$-subgroup of $G/N$, and $P \cap N$ is a Sylow $p$-subgroup of $N$.
\end{proposition}
\begin{proof}
By Proposition~\ref{prop:1.7}, $PN$ is a subgroup of $G$. The \nameref{thm:correspondence} applied to the natural map $G\to G/N$ gives
\[
    |G/N : PN/N| = |G : PN|.
\]
Since $P \leq PN \leq G$, the index $|G:PN|$ divides $|G:P|$. But $P$ is a Sylow $p$-subgroup of $G$, so
\[
    |G:P|=\frac{|G|}{|G|_p}
\]
is not divisible by $p$. Hence $p\nmid |G/N:PN/N|$. On the other hand, by the \nameref{thm:first-iso},
\[
    PN/N \cong P/(P\cap N),
\]
so $PN/N$ is a $p$-group. Thus $PN/N$ is a $p$-subgroup of $G/N$ whose index in $G/N$ is not divisible by $p$, and therefore $PN/N$ is a Sylow $p$-subgroup of $G/N$.\\
It remains to prove that $P\cap N$ is Sylow in $N$. Clearly $P\cap N$ is a $p$-subgroup. By Proposition~\ref{prop:3.12},
\[
    |PN|=\frac{|P||N|}{|P\cap N|},
\]
and hence
\[
    |N:P\cap N|=\frac{|N|}{|P\cap N|}=\frac{|PN|}{|P|}=|PN:P|.
\]
Since $P\leq PN\leq G$, the index $|PN:P|$ divides $|G:P|$, and again $p\nmid |G:P|$. Therefore $p\nmid |N:P\cap N|$. Hence $P\cap N$ is a $p$-subgroup of $N$ whose index in $N$ is not divisible by $p$, so $P\cap N$ is a Sylow $p$-subgroup of $N$.
\end{proof}

\subsubsection*{Applications: groups of order $pq$ and simplicity of $A_5$}
\begin{itemize}
    \item Sylow's theorem classifies groups of order $pq$ (distinct primes $p > q$).
    \item It also gives a clean count of $p$-elements via the orbit count of Sylow $p$-subgroups, used to prove $A_5$ is simple.
\end{itemize}

\begin{proposition}
\label{prop:7.3}
Let $p$ and $q$ be distinct primes, with $p > q$. If $p \not\equiv 1 \pmod{q}$, then any group of order $pq$ is isomorphic with $\Z_{pq}$. If $p \equiv 1 \pmod{q}$, then any abelian group of order $pq$ is isomorphic with $\Z_{pq}$, and there is exactly one isomorphism class of non-abelian groups of order $pq$.
\end{proposition}
\begin{proof}
Let $G$ be a group of order $pq$. Let $P$ be a Sylow $p$-subgroup of $G$, and let $Q$ be a Sylow $q$-subgroup of $G$. Then $|P|=p$ and $|Q|=q$, so $P\cong \Z_p$ and $Q\cong \Z_q$. Also, $P\cap Q=1$, and Proposition~\ref{prop:3.12} gives
\[
    |PQ|=\frac{|P||Q|}{|P\cap Q|}=pq=|G|,
\]
so $G=PQ$.\\
We first count the Sylow subgroups. By \nameref{thm:sylow} and Corollary~\ref{cor:7.1}, the number of Sylow $p$-subgroups divides $q$ and is congruent to $1$ modulo $p$. Since $p>q$, we have $q\not\equiv 1\pmod{p}$, so there is exactly one Sylow $p$-subgroup. Hence $P\trianglelefteq G$. Similarly, the number of Sylow $q$-subgroups divides $p$ and is congruent to $1$ modulo $q$. Therefore this number is either $1$ or $p$, and the second case can occur only when $p\equiv 1\pmod{q}$.\\
If $Q\trianglelefteq G$, then both $P$ and $Q$ are normal in $G$. Since $G=PQ$ and $P\cap Q=1$, Lemma~\ref{lem:2.7} gives
\[
    G\cong P\times Q\cong \Z_p\times \Z_q\cong \Z_{pq},
\]
where the final isomorphism follows from Lemma~\ref{lem:2.8}. Thus $G$ is cyclic whenever $Q$ is normal. This proves the first assertion when $p\not\equiv 1\pmod{q}$, and it also proves that every abelian group of order $pq$ is cyclic, since then $Q$ is automatically normal.\\
It remains to classify the non-abelian case, so assume $p\equiv 1\pmod{q}$ and that $G$ is non-abelian. Then $Q$ is not normal, while $P\trianglelefteq G$, $G=PQ$, and $P\cap Q=1$. Hence $G=P\rtimes Q$. Let
\[
    \varphi\colon Q\to \Aut(P)
\]
be the conjugation homomorphism. If $\ker\varphi\neq 1$, then $\ker\varphi=Q$, because $Q$ has prime order. In that case the action is trivial, so $G=P\times Q$, contradicting that $G$ is non-abelian. Therefore $\varphi$ is injective. Thus every non-abelian group of order $pq$ arises from a monomorphism
\[
    \Z_q\to \Aut(\Z_p).
\]
Conversely, any such monomorphism gives a non-abelian semidirect product $\Z_p\rtimes_{\varphi}\Z_q$ of order $pq$.\\
By Proposition~\ref{prop:2.2}, we have $\Aut(\Z_p)\cong \Z_{p-1}$. Since $q\mid (p-1)$, Theorem~\ref{thm:1.4} shows that $\Aut(\Z_p)$ has a unique subgroup $K$ of order $q$. Hence a monomorphism $\varphi\colon \Z_q\to \Aut(\Z_p)$ exists, and every such monomorphism has image $K$. If $\psi\colon \Z_q\to \Aut(\Z_p)$ is another monomorphism, then $\varphi(\Z_q)=K=\psi(\Z_q)$, so Proposition~\ref{prop:2.11} gives
\[
    \Z_p\rtimes_{\varphi}\Z_q\cong \Z_p\rtimes_{\psi}\Z_q.
\]
Therefore, when $p\equiv 1\pmod{q}$, there is exactly one isomorphism class of non-abelian groups of order $pq$.
\end{proof}

\begin{theorem}
\label{thm:7.4}
$A_5$ is simple.
\end{theorem}
\begin{proof}
We have
\[
    |A_5|=\frac{5!}{2}=60=2^2\cdot 3\cdot 5.
\]
We first count the elements of prime-power order that will be needed below. A $5$-cycle is even, and the number of $5$-cycles in $S_5$ is $4!=24$. Each Sylow $5$-subgroup of $A_5$ has order $5$ and hence contains $4$ non-identity elements, and two distinct subgroups of order $5$ intersect trivially. Thus $A_5$ has
\[
    \frac{24}{4}=6
\]
Sylow $5$-subgroups, and these account for all $24$ elements of order $5$. Similarly, the $3$-cycles are even, and there are
\[
    \binom{5}{3}\cdot 2=20
\]
of them. Each Sylow $3$-subgroup has two non-identity elements, so $A_5$ has $10$ Sylow $3$-subgroups, accounting for all $20$ elements of order $3$.\\
Now consider the Sylow $2$-subgroups. For $1\leq i\leq 5$, let $\{a,b,c,d\}$ be the complement of $\{i\}$ in $\{1,2,3,4,5\}$, and put
\[
    V_i=\{1,(ab)(cd),(ac)(bd),(ad)(bc)\}.
\]
Then $V_i$ is a subgroup of $A_5$ of order $4$, hence a Sylow $2$-subgroup. If $i\neq j$, then $V_i\cap V_j=1$, since a non-identity element of $V_i$ fixes exactly the point $i$. Moreover, for $\sigma\in A_5$ we have
\[
    \sigma V_i\sigma^{-1}=V_{\sigma(i)}.
\]
Since $A_5$ acts transitively on $\{1,2,3,4,5\}$, the five subgroups $V_1,\ldots,V_5$ are conjugate. By part (ii) of \nameref{thm:sylow}, these are all the Sylow $2$-subgroups of $A_5$. Thus $A_5$ has $5(4-1)=15$ elements of order $2$. Also, every element of order $2$ in $A_5$ has a distinct conjugate in $A_5$: for instance, if $x=(ab)(cd)$, with $i$ the remaining point, then conjugation by the $3$-cycle $(acb)$ sends $x$ to $(ac)(bd)\neq x$.\\
Now let $N\trianglelefteq A_5$ be a proper normal subgroup, and put $n=|N|$. By \nameref{thm:lagrange}, $n$ divides $60$, and since $N$ is proper we have $n\leq 30$. Suppose first that $5$ divides $n$. Then by \nameref{thm:cauchy}, $N$ contains an element of order $5$, and hence contains a Sylow $5$-subgroup of $A_5$. Since $N$ is normal in $A_5$, it contains every conjugate of this Sylow $5$-subgroup, so it contains all $24$ elements of order $5$. This forces $n=30$. But then $3$ divides $n$, so the same argument with Sylow $3$-subgroups shows that $N$ contains all $20$ elements of order $3$. These are disjoint from the $24$ elements of order $5$, and with the identity this gives more than $30$ elements in $N$, a contradiction. Therefore $5\nmid n$.\\
If $3$ divides $n$, then again by \nameref{thm:cauchy}, $N$ contains an element of order $3$, and normality forces $N$ to contain all $20$ elements of order $3$. But $5\nmid n$ and $n\mid 60$, so $n\leq 12$, a contradiction. Hence $3\nmid n$. Therefore $n$ is one of $1,2,4$. If $n=4$, then $N$ is a Sylow $2$-subgroup of $A_5$, but a Sylow $2$-subgroup has five conjugates, contradicting the normality of $N$. If $n=2$, then $N=\{1,x\}$ for some element $x$ of order $2$, but we have just observed that $x$ has a conjugate different from itself, again contradicting normality. Hence $n=1$.\\
Thus every proper normal subgroup of $A_5$ is trivial, and so $A_5$ is simple.
\end{proof}

\begin{theorem}
\label{thm:7.5}
Any simple group of order $60$ is isomorphic with $A_5$.
\end{theorem}
\begin{proof}
Let $G$ be a simple group of order $60$. We first observe that $G$ has no proper subgroup of index less than $5$. Indeed, if $H<G$ has index $r$, then the action of $G$ on $G/H$ by left multiplication gives a homomorphism
\[
    G\to S_r
\]
by Proposition~\ref{prop:3.1}. Since this action is transitive and $r>1$, it is non-trivial. The kernel is therefore a proper normal subgroup of the simple group $G$, and hence the kernel is trivial. Thus $G$ embeds in $S_r$. But if $r<5$, then
\[
    |S_r|=r!<60=|G|,
\]
which is impossible.\\
We next show that $G$ has a subgroup of index $5$. Suppose, for a contradiction, that it does not. Let $S$ be a Sylow $2$-subgroup of $G$. Then $|S|=4$. By Corollary~\ref{cor:7.1}, the number of Sylow $2$-subgroups divides $60/4=15$, and by Proposition~\ref{prop:3.14} this number is the index $|G:N_G(S)|$. This number is not $1$, since otherwise $S$ would be normal in $G$, contradicting simplicity. Since $G$ has no proper subgroup of index less than $5$ and, by assumption, none of index $5$, this number must be $15$.\\
Let $S_1$ and $S_2$ be distinct Sylow $2$-subgroups. Suppose that there is some $1\neq t\in S_1\cap S_2$. Since every group of order $4$ is abelian, both $S_1$ and $S_2$ lie in $C_G(t)$. Thus $|C_G(t)|>4$, and since $S_1\leq C_G(t)$, the order $|C_G(t)|$ is divisible by $4$. Hence
\[
    |G:C_G(t)|\leq 5.
\]
The first paragraph and our assumption that $G$ has no subgroup of index $5$ force $C_G(t)=G$. Thus $t\in Z(G)$, so $\langle t\rangle$ is a non-trivial proper normal subgroup of $G$, contradicting simplicity. Therefore distinct Sylow $2$-subgroups intersect trivially. Since there are $15$ Sylow $2$-subgroups, $G$ has
\[
    15(4-1)=45
\]
non-identity elements lying in Sylow $2$-subgroups.\\
Since $G$ is simple, its Sylow $5$-subgroup is not normal. By \nameref{thm:sylow} and Corollary~\ref{cor:7.1}, the number of Sylow $5$-subgroups divides $60/5=12$ and is congruent to $1$ modulo $5$, so this number must be $6$. These subgroups intersect trivially away from the identity, and hence contribute
\[
    6(5-1)=24
\]
elements of order $5$. This is impossible, since $45+24>60$. Therefore $G$ has a subgroup of index $5$.\\
Let $H$ be a subgroup of index $5$. The action of $G$ on $G/H$ gives a homomorphism $G\to S_5$. As above, simplicity forces this homomorphism to be injective, so we identify $G$ with its image in $S_5$. Then $|G|=60$, and hence $|S_5:G|=2$. By Proposition~\ref{prop:1.8}, $G\trianglelefteq S_5$.\\
We claim that $G=A_5$. Suppose not. Since $A_5\trianglelefteq S_5$, Proposition~\ref{prop:1.7} gives $GA_5\leq S_5$. The subgroup $GA_5$ contains $A_5$ properly, because $G\neq A_5$ and both groups have order $60$. Hence $|GA_5|>60$, so $GA_5=S_5$. By Proposition~\ref{prop:3.12},
\[
    |G\cap A_5|=\frac{|G||A_5|}{|GA_5|}
    =\frac{60\cdot 60}{120}
    =30.
\]
Therefore $G\cap A_5$ has index $2$ in $G$, so Proposition~\ref{prop:1.8} gives $G\cap A_5\trianglelefteq G$. This is a non-trivial proper normal subgroup of $G$, contradicting simplicity. Thus $G=A_5$, and the original group is isomorphic with $A_5$.
\end{proof}

\begin{corollary}
\label{cor:7.6}
$\mathrm{PSL}(2, 4) \cong \mathrm{PSL}(2, 5) \cong A_5$.
\end{corollary}
\begin{proof}
By Proposition~\ref{prop:6.5}, we have
\[
    |\mathrm{PSL}(2,4)|=4^3-4=60
\]
and
\[
    |\mathrm{PSL}(2,5)|=\frac{1}{2}(5^3-5)=60.
\]
By Theorem~\ref{thm:6.8}, both groups are simple, since the exceptional groups in degree $2$ occur only over fields of order $2$ and $3$. Therefore Theorem~\ref{thm:7.5} gives
\[
    \mathrm{PSL}(2,4)\cong A_5
    \qquad\text{and}\qquad
    \mathrm{PSL}(2,5)\cong A_5.
\]
Hence $\mathrm{PSL}(2,4)\cong \mathrm{PSL}(2,5)\cong A_5$.
\end{proof}

\noindent It is a relatively easy exercise to show that the only simple groups of order less than $60$ are the cyclic groups of prime order, and hence that $A_5$ is the non-abelian finite simple group having smallest order. The non-abelian finite simple group of next smallest order is $\mathrm{PSL}(2, 7)$, which has order $168$.

\subsection*{Exercises}

\noindent Throughout these exercises, $G$ is a finite group and $p$ is a prime divisor of $|G|$.

\begin{exercise}
\label{ex:7.1}
Let $H \leq G$ and let $P$ be a Sylow $p$-subgroup of $G$. Prove, without reference to \nameref{thm:sylow}, that there is some conjugate of $P$ whose intersection with $H$ is a Sylow $p$-subgroup of $H$.
\end{exercise}
\begin{solution}
Let $H$ act on the left cosets $G/P$ by left multiplication. Since $P$ is a Sylow $p$-subgroup of $G$, the index $|G:P|$ is not divisible by $p$. The $H$-orbits on $G/P$ partition $G/P$, so at least one such orbit has size not divisible by $p$. Choose a coset $gP$ in such an orbit.\\
The stabilizer of $gP$ in $H$ is
\[
    \operatorname{Stab}_H(gP)=\{h\in H: hgP=gP\}.
\]
Now
\[
    hgP=gP
    \quad\Longleftrightarrow\quad
    g^{-1}hg\in P
    \quad\Longleftrightarrow\quad
    h\in gPg^{-1},
\]
and hence
\[
    \operatorname{Stab}_H(gP)=H\cap gPg^{-1}.
\]
By Corollary~\ref{cor:3.5}, the size of the orbit of $gP$ is
\[
    |H:\operatorname{Stab}_H(gP)|=|H:H\cap gPg^{-1}|.
\]
This index is not divisible by $p$. Also, $H\cap gPg^{-1}$ is a $p$-subgroup of $H$, because it is contained in the $p$-group $gPg^{-1}$. Therefore $H\cap gPg^{-1}$ is a $p$-subgroup of $H$ whose index in $H$ is not divisible by $p$, so its order is the full $p$-part of $|H|$. Thus $H\cap gPg^{-1}$ is a Sylow $p$-subgroup of $H$.
\end{solution}

\begin{exercise}
\label{ex:7.2}
(cont.) Use Exercise~\ref{ex:7.1} to give an alternate proof of part (i) of \nameref{thm:sylow}.
\end{exercise}
\begin{solution}
Let $K$ be a finite group with $p\mid |K|$, and put $n=|K|$. By \nameref{thm:cayley}, we may regard $K$ as a subgroup of $S_n$. Embedding $S_n$ into $\GL(n,p)$ by permutation matrices, we may therefore regard
\[
    K\leq \GL(n,p).
\]
Let $U$ be the subgroup of upper uni-triangular matrices in $\GL(n,p)$. As computed at the beginning of this section,
\[
    |U|=p^{\frac{n(n-1)}{2}}=|\GL(n,p)|_p,
\]
so $U$ is a Sylow $p$-subgroup of $\GL(n,p)$ by definition. Applying Exercise~\ref{ex:7.1} to the subgroup $K\leq \GL(n,p)$ and the Sylow $p$-subgroup $U$, there is some $x\in \GL(n,p)$ such that
\[
    K\cap xUx^{-1}
\]
is a Sylow $p$-subgroup of $K$. Hence $K$ has a Sylow $p$-subgroup. Since $K$ was arbitrary, this proves part (i) of \nameref{thm:sylow}.
\end{solution}

\begin{exercise}
\label{ex:7.3}
Give an alternate proof of parts (ii) and (iii) of \nameref{thm:sylow} by considering the action of an arbitrary $p$-subgroup $Q$ of $G$ on the coset space $G/P$, where $P$ is a Sylow $p$-subgroup of $G$.
\end{exercise}
\begin{solution}
Let $Q$ be a $p$-subgroup of $G$, and let $Q$ act on $G/P$ by left multiplication:
\[
    q\cdot gP=qgP.
\]
Since $P$ is a Sylow $p$-subgroup of $G$, the index $|G:P|$ is not divisible by $p$. On the other hand, every orbit of the $p$-group $Q$ on $G/P$ has size a power of $p$, because its size is the index of a stabilizer in $Q$. Since the orbits partition $G/P$, at least one orbit must have size not divisible by $p$, and hence must have size $1$. Choose a fixed coset $gP$.\\
For every $q\in Q$, we have
\[
    qgP=gP.
\]
Equivalently,
\[
    g^{-1}qgP=P,
\]
so $g^{-1}qg\in P$, and therefore $q\in gPg^{-1}$. Thus
\[
    Q\leq gPg^{-1}.
\]
Since $gPg^{-1}$ is a conjugate of $P$, it is a Sylow $p$-subgroup of $G$. This proves part (iii): every $p$-subgroup of $G$ is contained in a Sylow $p$-subgroup.\\
Now let $Q$ itself be a Sylow $p$-subgroup of $G$. The containment just proved gives
\[
    Q\leq gPg^{-1}.
\]
But $Q$ and $gPg^{-1}$ have the same order, namely $|G|_p$. Hence $Q=gPg^{-1}$, so every Sylow $p$-subgroup of $G$ is conjugate to $P$. This proves part (ii).
\end{solution}

\begin{exercise}
\label{ex:7.4}
Prove the following generalization of part (iv) of \nameref{thm:sylow}: If $|G|$ is divisible by $p^b$, and $H \leq G$ has order $p^a$ where $a \leq b$, then the number of subgroups of $G$ that both contain $H$ and have order $p^b$ is congruent to $1$ modulo $p$.
\end{exercise}
\begin{solution}
We first record a one-step counting fact. Suppose that $K\leq G$ has order $p^c$ and that $p^{c+1}$ divides $|G|$. We claim that the number of subgroups $L\leq G$ such that
\[
    K\leq L,\qquad |L|=p^{c+1}
\]
is congruent to $1$ modulo $p$.\\
To see this, first note that if $M$ is any finite group whose order is divisible by $p$, then the number of subgroups of $M$ of order $p$ is congruent to $1$ modulo $p$. Indeed, let
\[
    X=\{(x_1,\ldots,x_p)\in M^p:x_1x_2\cdots x_p=1\}.
\]
Then $|X|=|M|^{p-1}$, so $p\mid |X|$. The cyclic group of order $p$ acts on $X$ by cyclically permuting the entries. The fixed points are exactly the tuples $(x,\ldots,x)$ with $x^p=1$. Hence the number of elements $x\in M$ satisfying $x^p=1$ is congruent to $0$ modulo $p$. If $r$ is the number of subgroups of $M$ of order $p$, then these elements are the identity together with the $p-1$ non-identity elements in each such subgroup, so
\[
    1+(p-1)r\equiv 0 \pmod{p}.
\]
Thus $r\equiv 1\pmod{p}$.\\
Now choose a Sylow $p$-subgroup $P$ of $G$ containing $K$, using part (iii) of \nameref{thm:sylow}. Since $p^{c+1}$ divides $|G|$, we have $K<P$. The $p$-group $K$ acts by left multiplication on the coset space $P/K$. The total number of cosets is divisible by $p$, and the coset $K$ is fixed. Since all non-trivial orbits have size divisible by $p$, there must be some other fixed coset $xK$. The condition that $xK$ is fixed says that $kxK=xK$ for every $k\in K$, or equivalently $x^{-1}kx\in K$ for every $k\in K$. Thus $x\in N_P(K)-K$, so $p$ divides $|N_G(K):K|$.\\
Subgroups $L\leq G$ with $K\leq L$ and $|L|=p^{c+1}$ are now in bijection with subgroups of $N_G(K)/K$ of order $p$. Indeed, if $K\leq L$ and $|L:K|=p$, then $K\trianglelefteq L$: the action of $L$ on the left cosets of $K$ gives a transitive homomorphism $L\to S_p$, whose image is a $p$-subgroup of $S_p$ and hence has order $p$, so the kernel has the same order as $K$ and is contained in $K$. Hence $L\leq N_G(K)$ and $L/K$ has order $p$. Conversely, each subgroup of $N_G(K)/K$ of order $p$ lifts to such an $L$. Since $p$ divides $|N_G(K)/K|$, the preceding paragraph applied to $M=N_G(K)/K$ proves the one-step claim.\\
We now prove the exercise by induction on $b-a$. If $b=a$, the only subgroup of order $p^b$ containing $H$ is $H$ itself, so the number is $1$. Assume $b>a$, and for $a\leq c\leq b$ put
\[
    \mathcal{A}_c=\{K\leq G:H\leq K,\ |K|=p^c\}.
\]
By induction, $|\mathcal{A}_{b-1}|\equiv 1\pmod{p}$. Count the pairs
\[
    (K,L)\quad\text{with}\quad K\in\mathcal{A}_{b-1},\ L\in\mathcal{A}_b,\ K\leq L.
\]
For each fixed $K\in\mathcal{A}_{b-1}$, the one-step claim shows that the number of possible $L$ is congruent to $1$ modulo $p$. Hence the total number of such pairs is congruent to $|\mathcal{A}_{b-1}|$, and therefore to $1$ modulo $p$.\\
On the other hand, fix $L\in\mathcal{A}_b$. The number of possible $K\in\mathcal{A}_{b-1}$ with $K\leq L$ is the number of subgroups of the $p$-group $L$ that contain $H$ and have order $p^{b-1}$. By the induction hypothesis applied inside $L$, this number is congruent to $1$ modulo $p$. Therefore the same total number of pairs is congruent to $|\mathcal{A}_b|$ modulo $p$. Combining the two counts gives
\[
    |\mathcal{A}_b|\equiv 1\pmod{p}.
\]
This is exactly the desired congruence.
\end{solution}

\begin{exercise}
\label{ex:7.5}
Show that if $P$ is a Sylow $p$-subgroup of $G$, then $N_G(P)$ is a self-normalizing subgroup of $G$, meaning that $N_G(N_G(P)) = N_G(P)$. (Any subgroup of $G$ that can be written as $N_G(P)$ for some non-trivial $p$-subgroup $P$ of $G$ is said to be a \emph{$p$-local subgroup} of $G$.)
\end{exercise}
\begin{solution}
Put $H=N_G(P)$. Since $P\leq H\leq G$ and $P$ has the full $p$-part of $|G|$, the subgroup $P$ also has the full $p$-part of $|H|$. Hence $P$ is a Sylow $p$-subgroup of $H$. Moreover, by the definition of normalizer, $P\trianglelefteq H$. Thus $P$ is the unique Sylow $p$-subgroup of $H$: if $Q$ is any Sylow $p$-subgroup of $H$, then part (ii) of \nameref{thm:sylow}, applied inside $H$, gives $Q=hPh^{-1}$ for some $h\in H$, and this equals $P$ because $P\trianglelefteq H$.\\
Now let $g\in N_G(H)$. Then $gHg^{-1}=H$. Since $P\leq H$, we have
\[
    gPg^{-1}\leq H.
\]
Also $gPg^{-1}$ has the same order as $P$, so it is a Sylow $p$-subgroup of $H$. By the uniqueness just proved, $gPg^{-1}=P$. Therefore $g\in N_G(P)=H$, and so
\[
    N_G(H)\leq H.
\]
The reverse inclusion is automatic, since every subgroup normalizes itself. Hence $H=N_G(H)$, which is exactly
\[
    N_G(N_G(P))=N_G(P).
\]
\end{solution}

\begin{exercise}
\label{ex:7.6}
Given a prime $p$, find an example of a finite group having exactly $1 + p$ Sylow $p$-subgroups. Can this be done for $1 + 2p$? $1 + 3p$?
\end{exercise}
\begin{solution}
First take
\[
    G=\GL(2,p)
\]
and let
\[
    U=\left\{
    \begin{pmatrix}
        1 & a\\
        0 & 1
    \end{pmatrix}
    :a\in\Z_p
    \right\}.
\]
Then $|U|=p$, while
\[
    |\GL(2,p)|=(p^2-1)(p^2-p)=p(p-1)^2(p+1),
\]
so $U$ is a Sylow $p$-subgroup of $G$. The normalizer of $U$ is the upper triangular subgroup
\[
    B=\left\{
    \begin{pmatrix}
        a & b\\
        0 & d
    \end{pmatrix}
    :a,d\in\Z_p^\times,\ b\in\Z_p
    \right\}.
\]
Indeed, $U$ fixes the line $\langle e_1\rangle$, and this is its unique fixed line in $\Z_p^2$; hence any element normalizing $U$ must preserve $\langle e_1\rangle$. Thus $N_G(U)=B$, and
\[
    |B|=p(p-1)^2.
\]
Therefore the number of Sylow $p$-subgroups of $\GL(2,p)$ is
\[
    |G:N_G(U)|=\frac{p(p-1)^2(p+1)}{p(p-1)^2}=p+1.
\]
So $1+p$ can be done for every prime $p$.\\
For the next two questions, there is a useful general construction. Suppose that $q$ is a prime power with $q\equiv 1\pmod p$. Let $F$ be the finite field of order $q$, and choose an element $\lambda\in F^\times$ of order $p$, which is possible because $F^\times$ is cyclic of order $q-1$. Let
\[
    G=F^+\rtimes \langle \lambda\rangle,
\]
where $\lambda$ acts on the additive group $F^+$ by multiplication. Then $|G|=qp$, and $A=\langle\lambda\rangle$ is a Sylow $p$-subgroup of $G$. We claim that $N_G(A)=A$. In the affine notation, $A$ consists of the maps $x\mapsto \lambda^i x$. If $t_b$ is the translation $x\mapsto x+b$, then
\[
    t_b(x\mapsto \lambda x)t_b^{-1}
\]
is the map
\[
    x\mapsto \lambda x+(1-\lambda)b.
\]
This lies again in $A$ only when $b=0$. Hence no non-trivial translation normalizes $A$, and so $N_G(A)=A$. Therefore
\[
    n_p(G)=|G:N_G(A)|=q.
\]
Thus $1+2p$ can be done whenever $1+2p$ is a prime power, and $1+3p$ can be done whenever $1+3p$ is a prime power. For example, $p=2$ gives $1+3p=7$, and $p=5$ gives $1+3p=16$. Also, for $p=3$, the group $A_5$ has exactly $10=1+3p$ Sylow $3$-subgroups, since $A_5$ has $20$ elements of order $3$ and each subgroup of order $3$ has two non-identity elements.\\
These constructions do not work for every prime. In fact, a theorem of M.~Hall, Jr.\ on possible Sylow numbers (``On the Number of Sylow Subgroups in a Finite Group'', \emph{J. Algebra} 7 (1967), Theorems 3.1--3.2) shows that $1+2p$ occurs exactly when $1+2p$ is a prime power, and that $1+3p$ occurs only for $p=2,3,5$. In particular, there is no finite group with $15$ Sylow $7$-subgroups, and no finite group with $1+3p$ Sylow $p$-subgroups for $p\geq 7$.
\end{solution}

\begin{exercise}
\label{ex:7.7}
Show that if $G$ has exactly $1 + kp$ Sylow $p$-subgroups for some $k \in \N$, then there is a subgroup of $S_{1+kp}$ having exactly $1 + kp$ Sylow $p$-subgroups.
\end{exercise}
\begin{solution}
Let $\Omega$ be the set of Sylow $p$-subgroups of $G$. By assumption,
\[
    |\Omega|=1+kp.
\]
The group $G$ acts on $\Omega$ by conjugation:
\[
    g\cdot P=gPg^{-1}.
\]
This gives a homomorphism
\[
    \theta\colon G\to \mathrm{Sym}(\Omega)\cong S_{1+kp}.
\]
Put $\overline{G}=\theta(G)$. Then $\overline{G}$ is a subgroup of $S_{1+kp}$. We claim that $\overline{G}$ has exactly $1+kp$ Sylow $p$-subgroups.\\
Let $K=\ker\theta$. Since $K\trianglelefteq G$, Proposition~\ref{prop:7.2} applied to the quotient $G/K$ shows that, if $P$ is a Sylow $p$-subgroup of $G$, then
\[
    \theta(P)\cong PK/K
\]
is a Sylow $p$-subgroup of $\overline{G}\cong G/K$. Moreover, every Sylow $p$-subgroup of $\overline{G}$ arises in this way from some Sylow $p$-subgroup of $G$.\\
It remains only to check that distinct Sylow $p$-subgroups of $G$ have distinct images in $\overline{G}$. Suppose that $P$ and $Q$ are Sylow $p$-subgroups of $G$ and that
\[
    \theta(P)=\theta(Q).
\]
Equivalently, $PK=QK$. Every element of $K$ fixes every point of $\Omega$, so in particular $K\leq N_G(P)$. Hence $PK$ is a subgroup of $G$ in which $P$ is normal. Since $P$ is a Sylow $p$-subgroup of $G$, it is also a Sylow $p$-subgroup of $PK$, and normality makes it the unique Sylow $p$-subgroup of $PK$. But $Q\leq QK=PK$, so $Q$ is also a Sylow $p$-subgroup of $PK$. Therefore $Q=P$.\\
Thus the map
\[
    P\mapsto \theta(P)
\]
gives a bijection from the Sylow $p$-subgroups of $G$ to the Sylow $p$-subgroups of $\overline{G}$. Hence $\overline{G}\leq S_{1+kp}$ has exactly $1+kp$ Sylow $p$-subgroups.
\end{solution}

\begin{exercise}
\label{ex:7.8}
Let $n \geq 6$, and assume by induction that $A_{n-1}$ is simple. Let $N$ be a non-trivial proper normal subgroup of $A_n$.
\begin{enumerate}
    \item[(a)] Show that $A_n = N \rtimes A_{n-1}$.
    \item[(b)] Let $N^{\#}$ be the set of non-identity elements of $N$, and consider the action of $A_{n-1}$ on $N^{\#}$ given by the conjugation homomorphism from $A_{n-1}$ to $\mathrm{Aut}(N)$. Show that $N^{\#}$ is isomorphic as an $A_{n-1}$-set to $\{1, \ldots, n-1\}$.
    \item[(c)] Show that $A_{n-1}$ acts triply transitively on $N^{\#}$. (An action of a group $G$ on a set $X$ is \emph{triply transitive} if for every pair of triples $(x_1, x_2, x_3)$ and $(y_1, y_2, y_3)$, where $x_i, y_i \in X$ for all $i$ and $x_i \neq x_j$ (resp., $y_i \neq y_j$) when $i \neq j$, there is some $g \in G$ such that $gx_i = y_i$ for all $i$.)
    \item[(d)] Derive a contradiction, and conclude that $A_n$ is simple.
\end{enumerate}
\end{exercise}
\begin{solution}
    \phantom{}\par
    \begin{enumerate}
        \item[(a)] Let $H=A_{n-1}$ be the stabilizer of $n$ in the natural action of $A_n$ on $\{1,\ldots,n\}$. First observe that this action is doubly transitive. Indeed, given two ordered pairs of distinct points, choose a permutation in $S_n$ sending the first pair to the second. If this permutation is odd, compose it on the right with a transposition of two points outside the first pair; this is possible because $n\geq 6$. The resulting permutation is even and still sends the first pair to the second. Thus Proposition~\ref{prop:3.8} shows that $H$ is a maximal subgroup of $A_n$.\\
        Now $N\cap H\trianglelefteq H$, because $N\trianglelefteq A_n$. By the induction hypothesis, $H\cong A_{n-1}$ is simple, so $N\cap H$ is either $1$ or $H$. If $N\cap H=H$, then $H\leq N$. Since $H$ is maximal and $N$ is proper, this would force $N=H$. But $H$ is not normal in $A_n$: its conjugates are the stabilizers of the other points. This contradicts $N\trianglelefteq A_n$. Hence
        \[
            N\cap H=1.
        \]
        Since $N$ is non-trivial, we have $N\nleq H$. Therefore $NH$ is a subgroup of $A_n$ properly containing $H$. By maximality of $H$, we get
        \[
            NH=A_n.
        \]
        Since $N\trianglelefteq A_n$ and $N\cap H=1$, this gives
        \[
            A_n=N\rtimes H=N\rtimes A_{n-1}.
        \]
        \item[(b)] From part (a), we have
        \[
            |N|=|A_n:H|=n.
        \]
        The subgroup $N$ acts on $\{1,\ldots,n\}$ by restriction of the natural action of $A_n$. The stabilizer in $N$ of the point $n$ is
        \[
            N\cap H=1,
        \]
        so the orbit of $n$ under $N$ has size $|N|=n$. Hence the map
        \[
            N\to \{1,\ldots,n\},\qquad x\mapsto x(n),
        \]
        is a bijection. It sends $1\in N$ to $n$, and therefore restricts to a bijection
        \[
            N^{\#}\to \{1,\ldots,n-1\}.
        \]
        This bijection is $H$-equivariant. Indeed, if $h\in H$ and $x\in N$, then $h$ fixes $n$, so
        \[
            (hxh^{-1})(n)=h x(h^{-1}(n))=h x(n).
        \]
        Thus conjugating $x$ by $h$ corresponds exactly to applying $h$ to the point $x(n)$. Therefore $N^{\#}$ is isomorphic as an $A_{n-1}$-set to $\{1,\ldots,n-1\}$.
        \item[(c)] By part (b), it is enough to prove that the natural action of $A_{n-1}$ on $\{1,\ldots,n-1\}$ is triply transitive. Let
        \[
            (a_1,a_2,a_3),\qquad (b_1,b_2,b_3)
        \]
        be ordered triples of distinct elements of $\{1,\ldots,n-1\}$. Choose a permutation $\sigma\in S_{n-1}$ such that $\sigma(a_i)=b_i$ for $i=1,2,3$. If $\sigma$ is even, then $\sigma\in A_{n-1}$ and we are done. If $\sigma$ is odd, then because $n-1\geq 5$, there are at least two points of $\{1,\ldots,n-1\}$ outside $\{a_1,a_2,a_3\}$. Let $\tau$ be the transposition of two such points. Then $\sigma\tau$ is even and still sends each $a_i$ to $b_i$. Hence $\sigma\tau\in A_{n-1}$ gives the required element. Therefore $A_{n-1}$ acts triply transitively on $N^{\#}$.
        \item[(d)] We now derive a contradiction from the triple transitivity of the conjugation action of $H=A_{n-1}$ on $N^{\#}$. Since $|N|=n\geq 6$, the set $N^{\#}$ has at least five elements. Choose $x\in N^{\#}$, and choose
        \[
            y\in N^{\#}-\{x,x^{-1}\}.
        \]
        Then $xy$ is non-identity and is distinct from both $x$ and $y$. Choose
        \[
            z\in N^{\#}-\{x,y,xy\}.
        \]
        By triple transitivity, there is some $h\in H$ such that
        \[
            hxh^{-1}=x,\qquad hyh^{-1}=y,\qquad h(xy)h^{-1}=z.
        \]
        But conjugation by $h$ is an automorphism of $N$, so
        \[
            h(xy)h^{-1}=(hxh^{-1})(hyh^{-1})=xy.
        \]
        This contradicts $z\neq xy$. Therefore no non-trivial proper normal subgroup $N$ of $A_n$ exists. Hence $A_n$ is simple.
    \end{enumerate}
\end{solution}

\noindent The remaining exercises develop a proof that if $G$ is a finite simple group with $|G| \leq 200$, then either $G \cong \Z_p$ for some prime $p$, or $G \cong A_5$, or $G \cong \mathrm{PSL}(2, 7)$. This was established by H\"older in 1892. We will need to assume two facts, which will be proved in later sections:
\begin{itemize}
    \item[A.] If $|G| = p^n$ where $n > 1$, then $G$ is not simple (Section~8).
    \item[B.] If $G$ is a $p$-group and $H < G$, then $H < N_G(H)$ (Section~11).
\end{itemize}

\begin{exercise}
\label{ex:7.9}
Show that $G$ is not simple whenever one of the following statements is true, where $p$ is an odd prime:
\begin{enumerate}
    \item[(a)] $|G| = 2^m p^n$, where $2^k \not\equiv 1 \pmod{p}$ for any $1 \leq k \leq m$.
    \item[(b)] $|G| = p^n q$, where $q \neq p$ is prime and $q \not\equiv 1 \pmod{p}$.
    \item[(c)] $|G| = ap$, where $p \nmid a$ and $kp + 1 \nmid a$ for any $k \in \N$.
\end{enumerate}
\end{exercise}
\begin{solution}
    We use the same Sylow-counting argument in each case. Let $n_p$ denote the number of Sylow $p$-subgroups of $G$. By \nameref{thm:sylow} and Corollary~\ref{cor:7.1},
    \[
        n_p\equiv 1\pmod p
        \quad\text{and}\quad
        n_p \mid |G:P|,
    \]
    where $P$ is a Sylow $p$-subgroup of $G$. If $n_p=1$, then $P$ is the unique Sylow $p$-subgroup of $G$, and hence $P\trianglelefteq G$. This gives a non-trivial proper normal subgroup of $G$, so $G$ is not simple.
    \phantom{}\par
    \begin{enumerate}
        \item[(a)] Suppose $|G|=2^m p^n$. Then a Sylow $p$-subgroup has order $p^n$, so
        \[
            n_p \mid 2^m.
        \]
        Hence $n_p=2^k$ for some $0\leq k\leq m$. If $k\geq 1$, then the Sylow congruence gives
        \[
            2^k=n_p\equiv 1\pmod p,
        \]
        contradicting the hypothesis. Therefore $k=0$, and so $n_p=1$. Hence the Sylow $p$-subgroup is normal, and $G$ is not simple.
        \item[(b)] Suppose $|G|=p^nq$, where $q\neq p$ is prime. Then
        \[
            n_p\mid q
            \quad\text{and}\quad
            n_p\equiv 1\pmod p.
        \]
        Since $q$ is prime, we have $n_p=1$ or $n_p=q$. The second possibility would force
        \[
            q\equiv 1\pmod p,
        \]
        contrary to the hypothesis. Thus $n_p=1$, so the Sylow $p$-subgroup is normal and $G$ is not simple.
        \item[(c)] Suppose $|G|=ap$, where $p\nmid a$. Then a Sylow $p$-subgroup has order $p$, so
        \[
            n_p\mid a
            \quad\text{and}\quad
            n_p\equiv 1\pmod p.
        \]
        If $n_p>1$, then $n_p=kp+1$ for some $k\in\N$. But this is impossible, because the hypothesis says that no such number $kp+1$ divides $a$. Hence $n_p=1$. Therefore the Sylow $p$-subgroup is normal, and $G$ is not simple.
    \end{enumerate}
\end{solution}

\begin{exercise}
\label{ex:7.10}
Show that $G$ is not simple whenever one of the following statments is true:
\begin{enumerate}
    \item[(a)] $|G| = p^a(p + 1)$, where $a > 1$.
    \item[(b)] $|G| = p^a(p + 3)$, where $a > 3$ if $p = 2$ and $a > 1$ if $p > 3$.
    \item[(c)] $|G| = p^a(p^2 - 1)$, where $a > 1$ and $p$ is odd.
\end{enumerate}
\end{exercise}
\begin{solution}
    We again write $n_p$ for the number of Sylow $p$-subgroups of $G$. If $n_p=1$, then the Sylow $p$-subgroup is normal, so $G$ is not simple. We may therefore assume, in each case, that $n_p>1$ and derive a contradiction.
    \phantom{}\par
    \begin{enumerate}
        \item[(a)] Let $|G|=p^a(p+1)$ with $a>1$, and let $P$ be a Sylow $p$-subgroup of $G$. By \nameref{thm:sylow} and Corollary~\ref{cor:7.1},
        \[
            n_p\equiv 1\pmod p
            \quad\text{and}\quad
            n_p\mid p+1.
        \]
        Thus $n_p=1$ or $n_p=p+1$. If $G$ were simple, then $n_p\neq 1$, so $n_p=p+1$. The conjugation action of $G$ on its Sylow $p$-subgroups gives a homomorphism
        \[
            G\to S_{p+1}.
        \]
        Its kernel is normal in $G$, and the action is non-trivial because there is more than one Sylow $p$-subgroup. Hence, if $G$ is simple, the kernel is trivial and $G$ embeds in $S_{p+1}$. Therefore
        \[
            |G|\mid (p+1)!.
        \]
        But $p^2\mid |G|$ because $a>1$, whereas $p^2\nmid (p+1)!$, since among the integers $1,\ldots,p+1$ there is only one multiple of $p$. This contradiction shows that $G$ is not simple.
        \item[(b)] First suppose $p=2$. Then $|G|=2^a\cdot 5$ with $a>3$. The number $n_2$ of Sylow $2$-subgroups satisfies
        \[
            n_2\mid 5
            \quad\text{and}\quad
            n_2\equiv 1\pmod 2.
        \]
        Thus $n_2=1$ or $n_2=5$. If $G$ were simple, then $n_2=5$, and the conjugation action on the five Sylow $2$-subgroups would embed $G$ in $S_5$. Hence $2^a\cdot 5$ would divide $5!=120=2^3\cdot 3\cdot 5$, contradicting $a>3$. Therefore $G$ is not simple in this case.\\
        Now suppose $p>3$. Then $p\nmid p+3$, so a Sylow $p$-subgroup has order $p^a$, and
        \[
            n_p\mid p+3
            \quad\text{and}\quad
            n_p\equiv 1\pmod p.
        \]
        Since $n_p\leq p+3$, the only values congruent to $1$ modulo $p$ are $1$ and $p+1$. But $p+1\nmid p+3$. Hence $n_p=1$, so the Sylow $p$-subgroup is normal and $G$ is not simple.
        \item[(c)] Let $|G|=p^a(p^2-1)$, where $p$ is odd and $a>1$. Then
        \[
            n_p\mid p^2-1
            \quad\text{and}\quad
            n_p\equiv 1\pmod p.
        \]
        Write $n_p=kp+1$. Since $n_p\mid p^2-1$, we have $0\leq k\leq p-1$. Also $kp\equiv -1\pmod{n_p}$, so
        \[
            k^2p^2\equiv 1\pmod{n_p}.
        \]
        As $n_p\mid p^2-1$, it follows that
        \[
            n_p\mid k^2(p^2-1).
        \]
        Combining this with the congruence above gives
        \[
            n_p\mid k^2-1.
        \]
        If $k\geq 2$, then $n_p=kp+1>k^2-1$, because $k\leq p-1$. This is impossible. Hence $k=0$ or $k=1$, and therefore
        \[
            n_p=1
            \quad\text{or}\quad
            n_p=p+1.
        \]
        If $G$ were simple, then $n_p\neq 1$, so $n_p=p+1$. As in part (a), the conjugation action on the Sylow $p$-subgroups would embed $G$ in $S_{p+1}$. Thus $|G|\mid (p+1)!$. But $p^2\mid |G|$ and $p^2\nmid (p+1)!$, contradicting $a>1$. Therefore $G$ is not simple.
    \end{enumerate}
\end{solution}

\begin{exercise}
\label{ex:7.11}
Show that if $|G| \in \{30, 56, 105, 132\}$, then $G$ is not simple.
\end{exercise}
\begin{solution}
    We argue by contradiction in each case, assuming that $G$ is simple.
    \phantom{}\par
    \begin{enumerate}
        \item If $|G|=30=2\cdot 3\cdot 5$, then
        \[
            n_5\mid 6,\qquad n_5\equiv 1\pmod 5,
        \]
        so $n_5=1$ or $6$. Simplicity forces $n_5=6$. These six Sylow $5$-subgroups are distinct subgroups of order $5$, and hence intersect pairwise only in the identity. They therefore contribute $6(5-1)=24$ non-identity elements. Similarly,
        \[
            n_3\mid 10,\qquad n_3\equiv 1\pmod 3,
        \]
        so $n_3=1$ or $10$. Simplicity forces $n_3=10$, and these Sylow $3$-subgroups contribute $10(3-1)=20$ non-identity elements. Elements of order $3$ and elements of order $5$ are disjoint, so this gives at least $24+20=44$ non-identity elements, impossible in a group of order $30$.
        \item If $|G|=56=2^3\cdot 7$, then
        \[
            n_7\mid 8,\qquad n_7\equiv 1\pmod 7,
        \]
        so $n_7=1$ or $8$. Simplicity forces $n_7=8$. These eight Sylow $7$-subgroups contribute $8(7-1)=48$ non-identity elements. Hence only $56-1-48=7$ non-identity elements remain outside the Sylow $7$-subgroups. On the other hand,
        \[
            n_2\mid 7,\qquad n_2\equiv 1\pmod 2,
        \]
        so $n_2=1$ or $7$. Simplicity forces $n_2=7$. Each Sylow $2$-subgroup has order $8$, so its seven non-identity elements must be exactly the seven remaining non-identity elements. Thus all Sylow $2$-subgroups are equal, contradicting $n_2=7$.
        \item If $|G|=105=3\cdot 5\cdot 7$, then
        \[
            n_7\mid 15,\qquad n_7\equiv 1\pmod 7,
        \]
        so $n_7=1$ or $15$. Simplicity forces $n_7=15$, giving $15(7-1)=90$ non-identity elements of order $7$. Also
        \[
            n_5\mid 21,\qquad n_5\equiv 1\pmod 5,
        \]
        so $n_5=1$ or $21$. Simplicity forces $n_5=21$, giving $21(5-1)=84$ non-identity elements of order $5$. These are disjoint from the elements of order $7$, giving more elements than $G$ has. This is impossible.
        \item If $|G|=132=2^2\cdot 3\cdot 11$, then
        \[
            n_{11}\mid 12,\qquad n_{11}\equiv 1\pmod {11},
        \]
        so $n_{11}=1$ or $12$. Simplicity forces $n_{11}=12$, and these Sylow $11$-subgroups contribute $12(11-1)=120$ non-identity elements. Now
        \[
            n_3\mid 44,\qquad n_3\equiv 1\pmod 3,
        \]
        so $n_3\in\{1,4,22\}$. Simplicity gives $n_3\neq 1$. If $n_3=22$, then the Sylow $3$-subgroups contribute $22(3-1)=44$ non-identity elements, disjoint from the elements of order $11$, which is impossible. Hence $n_3=4$. The conjugation action of $G$ on its four Sylow $3$-subgroups gives a non-trivial homomorphism
        \[
            G\to S_4.
        \]
        Since $G$ is simple, the kernel must be trivial, so $G$ embeds in $S_4$. This is impossible because $132\nmid 24=|S_4|$.
    \end{enumerate}
\end{solution}

\begin{exercise}
\label{ex:7.12}
Show that $G$ is not simple in the following cases:
\begin{enumerate}
    \item[(a)] $|G| = 90$. (\textsc{Hint}: Show that if $Q$ is a Sylow $3$-subgroup of $G$ and $1 \neq x \in Q$, then $C_G(x) = Q$.)
    \item[(b)] $|G| = 112$. (\textsc{Hint}: Use fact B above to show that any two Sylow $2$-subgroups of $G$ intersect trivially.)
    \item[(c)] $|G| = 120$. (\textsc{Hint}: Show that $G$ imbeds in $S_6$, and contradict Exercise~\ref{ex:3.8}.)
    \item[(d)] $|G| = 144$. (\textsc{Hint}: Argue as in part (i) above.)
    \item[(e)] $|G| = 180$. (\textsc{Hint}: Show that if $Q$ is a Sylow $3$-subgroup of $G$ and $1 \neq x \in Q$, then $|C_G(x) : Q| \leq 2$.)
\end{enumerate}
\end{exercise}
\begin{solution}
    We argue by contradiction in each case, assuming that $G$ is simple.
    \phantom{}\par
    \begin{enumerate}
        \item[(a)] Let $|G|=90=2\cdot 3^2\cdot 5$. Then
        \[
            n_3\mid 10,\qquad n_3\equiv 1\pmod 3,
        \]
        so simplicity forces $n_3=10$. Also
        \[
            n_5\mid 18,\qquad n_5\equiv 1\pmod 5,
        \]
        so simplicity forces $n_5=6$. Let $Q$ be a Sylow $3$-subgroup and let $1\neq x\in Q$. Since $|Q|=9$, the group $Q$ is abelian, so $Q\leq C_G(x)$. Hence $|C_G(x)|$ is divisible by $9$. If $|C_G(x)|=90$, then $\langle x\rangle$ is a non-trivial proper normal subgroup of $G$. If $|C_G(x)|=45$, then $C_G(x)$ has index $2$ and is normal in $G$. If $|C_G(x)|=18$, then the action of $G$ on the five cosets of $C_G(x)$ gives a non-trivial homomorphism $G\to S_5$, which must be injective by simplicity; this is impossible because $90\nmid 5!$. Therefore $|C_G(x)|=9$, and so
        \[
            C_G(x)=Q.
        \]
        It follows that distinct Sylow $3$-subgroups intersect trivially: if two of them contain the same non-identity element $x$, then both are contained in $C_G(x)$ and hence are equal. Thus the Sylow $3$-subgroups contribute $10(9-1)=80$ non-identity elements. The Sylow $5$-subgroups contribute $6(5-1)=24$ further non-identity elements, disjoint from these. This is impossible in a group of order $90$.
        \item[(b)] Let $|G|=112=2^4\cdot 7$. Then
        \[
            n_2\mid 7,\qquad n_2\equiv 1\pmod 2,
        \]
        and
        \[
            n_7\mid 16,\qquad n_7\equiv 1\pmod 7.
        \]
        Simplicity forces $n_2=7$ and $n_7=8$. We first show that distinct Sylow $2$-subgroups intersect trivially. Suppose otherwise, and choose distinct Sylow $2$-subgroups $S$ and $T$ such that
        \[
            D=S\cap T\neq 1
        \]
        has maximal possible order among intersections of distinct Sylow $2$-subgroups. By fact B, applied inside the $2$-groups $S$ and $T$, we have
        \[
            D<N_S(D)
            \quad\text{and}\quad
            D<N_T(D).
        \]
        Put $N=N_G(D)$. We claim that $N$ has more than one Sylow $2$-subgroup. If $U$ were the unique Sylow $2$-subgroup of $N$, then $U$ would contain both $N_S(D)$ and $N_T(D)$. Let $R$ be a Sylow $2$-subgroup of $G$ containing $U$. Since
        \[
            N_S(D)\leq S\cap R
            \quad\text{and}\quad
            N_S(D)>D,
        \]
        the maximality of $D$ forces $R=S$. The same argument with $T$ forces $R=T$, a contradiction. Hence $N$ has more than one Sylow $2$-subgroup. The number of Sylow $2$-subgroups of $N$ is odd and greater than $1$, so $7$ divides $|N|$. Also $N$ contains the $2$-subgroup $N_S(D)$, which has order at least $4$. Hence $28$ divides $|N|$. Thus $|G:N|$ is $1$, $2$, or $4$. If $|G:N|=1$, then $D\trianglelefteq G$; if $|G:N|=2$, then $N\trianglelefteq G$; and if $|G:N|=4$, then the coset action gives an embedding $G\hookrightarrow S_4$, impossible because $112\nmid 24$. This proves that distinct Sylow $2$-subgroups intersect trivially.\\
        Therefore the seven Sylow $2$-subgroups contribute $7(16-1)=105$ non-identity elements. The eight Sylow $7$-subgroups contribute $8(7-1)=48$ non-identity elements, disjoint from the elements in the Sylow $2$-subgroups. This is impossible in a group of order $112$.
        \item[(c)] Let $|G|=120=2^3\cdot 3\cdot 5$. Then
        \[
            n_5\mid 24,\qquad n_5\equiv 1\pmod 5,
        \]
        so simplicity forces $n_5=6$. The conjugation action of $G$ on its six Sylow $5$-subgroups gives a non-trivial homomorphism
        \[
            G\to S_6.
        \]
        Since $G$ is simple, this homomorphism is injective. Let $H$ be its image. The sign map on $H$ must be trivial, since otherwise $H\cap A_6$ would be a proper normal subgroup of the simple group $H$. Hence $H\leq A_6$. But then
        \[
            |A_6:H|=\frac{360}{120}=3.
        \]
        By Exercise~\ref{ex:7.8}, the group $A_6$ is simple. Its action on the three cosets of $H$ would therefore give an embedding $A_6\hookrightarrow S_3$, impossible because $360\nmid 6$. Thus $G$ is not simple.
        \item[(d)] Let $|G|=144=2^4\cdot 3^2$. Then
        \[
            n_3\mid 16,\qquad n_3\equiv 1\pmod 3,
        \]
        so $n_3\in\{1,4,16\}$. Simplicity gives $n_3\neq 1$. If $n_3=4$, then the conjugation action on the four Sylow $3$-subgroups embeds $G$ in $S_4$, impossible because $144\nmid 24$. Hence $n_3=16$.\\
        We claim that distinct Sylow $3$-subgroups intersect trivially. Suppose $Q$ and $R$ are distinct Sylow $3$-subgroups with $D=Q\cap R\neq 1$. Since $|Q|=|R|=9$, both $Q$ and $R$ are abelian, and hence both lie in $N_G(D)$. Thus $N_G(D)$ has more than one Sylow $3$-subgroup. The number of its Sylow $3$-subgroups divides $16$ and is congruent to $1$ modulo $3$, so it is at least $4$. Hence $36$ divides $|N_G(D)|$. Therefore $|G:N_G(D)|$ is $1$, $2$, or $4$. These cases are impossible respectively because they would make $D$ normal in $G$, make $N_G(D)$ normal in $G$, or embed $G$ in $S_4$. Thus distinct Sylow $3$-subgroups intersect trivially.\\
        The sixteen Sylow $3$-subgroups therefore contribute $16(9-1)=128$ non-identity elements. Only $144-1-128=15$ non-identity elements remain. But a Sylow $2$-subgroup has order $16$, and hence has exactly $15$ non-identity elements, all outside the Sylow $3$-subgroups. Therefore every Sylow $2$-subgroup is the same subgroup, making it normal. This contradicts simplicity.
        \item[(e)] Let $|G|=180=2^2\cdot 3^2\cdot 5$. Let $Q$ be a Sylow $3$-subgroup and let $1\neq x\in Q$. Since $|Q|=9$, the group $Q$ is abelian, so $Q\leq C_G(x)$. Thus $|C_G(x)|$ is divisible by $9$. The possibilities are $9$, $18$, $36$, $45$, $90$, and $180$. If $|C_G(x)|=180$, then $\langle x\rangle\trianglelefteq G$. If $|C_G(x)|=90$, then $C_G(x)$ has index $2$ and is normal in $G$. If $|C_G(x)|=45$, then the coset action gives an embedding $G\hookrightarrow S_4$, impossible because $180\nmid 24$. If $|C_G(x)|=36$, then the coset action gives an embedding $G\hookrightarrow S_5$, impossible because $180\nmid 120$. Hence
        \[
            |C_G(x):Q|\leq 2.
        \]
        In particular, $Q$ is the unique Sylow $3$-subgroup of $C_G(x)$. Therefore distinct Sylow $3$-subgroups intersect trivially. Now
        \[
            n_3\mid 20,\qquad n_3\equiv 1\pmod 3,
        \]
        so $n_3\in\{1,4,10\}$, and simplicity gives $n_3\neq 1$. If $n_3=4$, then the non-identity elements in the Sylow $3$-subgroups number $4(9-1)=32$. But each conjugacy class of such an element has size $|G:C_G(x)|$, which is either $20$ or $10$, contradicting that $32$ is not a sum of numbers of this form. Hence $n_3=10$, and there are $10(9-1)=80$ non-identity elements lying in Sylow $3$-subgroups.\\
        Finally,
        \[
            n_5\mid 36,\qquad n_5\equiv 1\pmod 5,
        \]
        so $n_5\in\{1,6,36\}$. Simplicity gives $n_5\neq 1$. If $n_5=36$, then the Sylow $5$-subgroups contribute $36(5-1)=144$ non-identity elements, disjoint from the $80$ elements already counted in Sylow $3$-subgroups, impossible. Hence $n_5=6$. The conjugation action on the six Sylow $5$-subgroups embeds $G$ in $S_6$. Its image lies in $A_6$ by the sign-map argument used in part (c), and has order $180$, hence index $2$ in $A_6$. This contradicts the simplicity of $A_6$.
    \end{enumerate}
\end{solution}

\noindent From fact A above and Exercises 9--12, we can conclude that if $G$ is a simple group with $|G| \leq 200$, then $|G|$ either is prime or is equal to either $60$ or $168$. (Verify this.) In light of Theorem~\ref{thm:7.5} and Theorem~\ref{thm:6.8}, all we need now show is that any two simple groups of order $168 = 2^3 \cdot 3 \cdot 7$ are isomorphic. We accomplish this in Exercises 13--17 below by showing that any simple group of order $168$ is isomorphic with a specific subgroup of $A_8$. (Compare with Exercise~\ref{ex:6.9}.) Let $G$ be a simple group of order $168$, let $P = \langle x \rangle$ be a Sylow $7$-subgroup of $G$, and let $H = N_G(P)$.

\begin{exercise}
\label{ex:7.13}
Show that $H = \{g \in G \mid gxg^{-1} \in P\}$ is a non-abelian group of order $21$ which is generated by $x$ and $y$, where $y$ has order $3$ and $yxy^{-1} = x^4$.
\end{exercise}
\begin{solution}
    Since $P=\langle x\rangle$, an element $g\in G$ normalizes $P$ if and only if
    \[
        gPg^{-1}=P.
    \]
    If $g$ normalizes $P$, then certainly $gxg^{-1}\in P$. Conversely, if $gxg^{-1}\in P$, then $gxg^{-1}$ has order $7$, so it is a non-identity element of the cyclic group $P$ of order $7$. Hence it generates $P$, and therefore
    \[
        gPg^{-1}=\langle gxg^{-1}\rangle=P.
    \]
    Thus
    \[
        H=N_G(P)=\{g\in G\mid gxg^{-1}\in P\}.
    \]
    Next,
    \[
        n_7\mid 24,\qquad n_7\equiv 1\pmod 7,
    \]
    so $n_7=1$ or $8$. Since $G$ is simple, $n_7\neq 1$, and hence $n_7=8$. The Sylow $7$-subgroups are the conjugates of $P$, so
    \[
        8=n_7=|G:N_G(P)|=|G:H|.
    \]
    Therefore
    \[
        |H|=\frac{168}{8}=21.
    \]
    We now show that $H$ is non-abelian. Suppose, for a contradiction, that $H$ is abelian. By \nameref{thm:cauchy}, $H$ contains a subgroup $Y$ of order $3$. Since $H$ is abelian, $Y\trianglelefteq H$, so $H\leq N_G(Y)$. Thus $|N_G(Y)|$ is divisible by $21$, and hence
    \[
        |N_G(Y)|\in\{21,42,84,168\}.
    \]
    If $|N_G(Y)|=168$, then $Y\trianglelefteq G$. If $|N_G(Y)|=84$, then $N_G(Y)$ has index $2$ in $G$ and is normal. If $|N_G(Y)|=42$, then the action of $G$ on the four cosets of $N_G(Y)$ gives an embedding $G\hookrightarrow S_4$, impossible because $168\nmid 24$. Finally, if $|N_G(Y)|=21$, then the number of Sylow $3$-subgroups of $G$ is
    \[
        |G:N_G(Y)|=8,
    \]
    contradicting the congruence $n_3\equiv 1\pmod 3$. Hence $H$ is non-abelian.\\
    Choose an element $y\in H$ of order $3$, again by \nameref{thm:cauchy}. Since $P\trianglelefteq H$, the subgroup $P\langle y\rangle$ is a subgroup of $H$ of order
    \[
        \frac{|P|\,|\langle y\rangle|}{|P\cap \langle y\rangle|}=7\cdot 3=21.
    \]
    Thus
    \[
        H=P\langle y\rangle=\langle x,y\rangle.
    \]
    Since $y\in H=N_G(P)$, conjugation by $y$ restricts to an automorphism of $P$. Therefore
    \[
        yxy^{-1}=x^a
    \]
    for some $a\in\{1,2,3,4,5,6\}$. Because $y$ has order $3$, this automorphism has order dividing $3$, so
    \[
        a^3\equiv 1\pmod 7.
    \]
    Since $H$ is non-abelian, the action is not trivial, so $a\neq 1$. The non-trivial solutions are $a=2$ and $a=4$. If necessary, replacing $y$ by $y^2$ changes $a$ to $a^2$, so we may choose $y$ so that
    \[
        yxy^{-1}=x^4.
    \]
\end{solution}

\begin{exercise}
\label{ex:7.14}
(cont.) Show that $N_G(\langle y \rangle)$ is isomorphic with $S_3$ and is generated by $y$ and $z$, where $z$ has order $2$ and $yzy = z$.
\end{exercise}
\begin{solution}
    Put $Y=\langle y\rangle$ and $N=N_G(Y)$. We first compute the number of Sylow $3$-subgroups of $G$. Inside $H$, the subgroup $P=\langle x\rangle$ is normal. If $Y$ were also normal in $H$, then $P$ and $Y$ would be normal subgroups of coprime orders, so they would commute elementwise and $H=PY$ would be abelian. This contradicts Exercise~\ref{ex:7.13}. Hence $Y$ is not normal in $H$, and the number of Sylow $3$-subgroups of $H$ is $7$.\\
    Now the Sylow congruence in $G$ gives
    \[
        n_3\mid 56,\qquad n_3\equiv 1\pmod 3,
    \]
    so $n_3\in\{1,4,7,28\}$. Simplicity gives $n_3\neq 1$. If $n_3=4$, then the conjugation action on the four Sylow $3$-subgroups gives an embedding $G\hookrightarrow S_4$, impossible because $168\nmid 24$. We also cannot have $n_3=7$. Indeed, every conjugate of $H$ contains seven Sylow $3$-subgroups, so if $G$ had only seven Sylow $3$-subgroups, then every conjugate of $H$ would contain all of them. But there are eight conjugates of $H$, one for each Sylow $7$-subgroup, and two distinct conjugates of $H$ would then have intersection containing the identity and all fourteen non-identity elements in the Sylow $3$-subgroups. This is impossible, since the intersection of two subgroups of order $21$ has order dividing $21$. Therefore
    \[
        n_3=28.
    \]
    Since the Sylow $3$-subgroups are the conjugates of $Y$, we get
    \[
        |G:N_G(Y)|=28.
    \]
    Hence
    \[
        |N|=|N_G(Y)|=\frac{168}{28}=6.
    \]
    We claim that $N$ is not cyclic. Suppose otherwise. Then, for each Sylow $3$-subgroup $Y'$ of $G$, the group $N_G(Y')$ is cyclic of order $6$, and hence contains exactly two elements of order $6$. Distinct Sylow $3$-subgroups give disjoint such elements of order $6$, because an element of order $6$ determines its unique subgroup of order $3$. Thus $G$ would contain
    \[
        28\cdot 2=56
    \]
    elements of order $6$. The Sylow $3$-subgroups contribute $28\cdot 2=56$ elements of order $3$, and the Sylow $7$-subgroups contribute $8\cdot 6=48$ elements of order $7$. Altogether this accounts for
    \[
        56+56+48=160
    \]
    non-identity elements. Only seven non-identity elements remain. But a Sylow $2$-subgroup has order $8$, so its seven non-identity elements must be exactly these seven remaining elements. Hence the Sylow $2$-subgroup is unique and normal, contradicting the simplicity of $G$. Therefore $N$ is non-cyclic.\\
    A non-cyclic group of order $6$ is isomorphic with $S_3$. By \nameref{thm:cauchy}, choose an element $z\in N$ of order $2$. Then $N=\langle y,z\rangle$, since this subgroup has order divisible by both $2$ and $3$ and is contained in the group $N$ of order $6$. Since $z\in N_G(Y)$, conjugation by $z$ sends $y$ either to $y$ or to $y^{-1}$. It cannot send $y$ to itself, since then $y$ and $z$ would commute and $N=\langle y,z\rangle$ would be cyclic. Therefore
    \[
        zyz^{-1}=y^{-1}.
    \]
    As $z$ has order $2$, this is equivalent to
    \[
        yzy=z.
    \]
    Thus $N_G(\langle y\rangle)=\langle y,z\rangle\cong S_3$, with $z$ of order $2$ and $yzy=z$.
\end{solution}

\begin{exercise}
\label{ex:7.15}
(cont.) Show that $G/H = \{H, zH, xzH, \ldots, x^6 zH\}$. Conclude that $G = \langle x, y, z \rangle$.
\end{exercise}
\begin{solution}
    We know from Exercise~\ref{ex:7.13} that $|H|=21$, so
    \[
        |G:H|=\frac{168}{21}=8.
    \]
    Thus it is enough to show that the eight displayed left cosets are distinct. First, $zH\neq H$, because otherwise $z\in H$, impossible since $z$ has order $2$ while $|H|=21$. Similarly, $x^i zH\neq H$ for every $0\leq i\leq 6$, because $x^i\in H$ and $x^iz\in H$ would force $z\in H$.\\
    Now suppose that
    \[
        x^i zH=x^j zH
    \]
    for some $0\leq i,j\leq 6$. Then
    \[
        z x^{i-j}z\in H.
    \]
    If $i\neq j$, then $x^{i-j}$ is a non-identity element of $P$, and hence generates $P$. Since $z x^{i-j}z$ is a non-identity element of order $7$ lying in $H$, it lies in the unique Sylow $7$-subgroup $P$ of $H$. Therefore
    \[
        zPz^{-1}=P,
    \]
    so $z\in N_G(P)=H$, again impossible. Hence $i=j$. The eight cosets
    \[
        H,\ zH,\ xzH,\ldots,x^6zH
    \]
    are therefore distinct, and since there are exactly eight left cosets of $H$ in $G$, they are all of them.\\
    Finally, Exercise~\ref{ex:7.13} gives $H=\langle x,y\rangle$. Let $K=\langle x,y,z\rangle$. Then $H\leq K$, and each coset representative in the displayed list lies in $K$. Hence every element of $G$ lies in $K$, so
    \[
        G=K=\langle x,y,z\rangle.
    \]
\end{solution}

\noindent Let $\varphi \colon G \to S_8$ be the monomorphism corresponding to the action of $G$ on the coset space $G/H$. Observe that $\varphi(G) = \varphi(G)' \leq S_8' \leq A_8$. View $S_8$ as the group of permutations of the set $\{\infty, 0, 1, 2, 3, 4, 5, 6\}$; for each $0 \leq i \leq 6$, associate $i$ with the coset $x^i z H$, and associate $\infty$ with the coset $H$.

\begin{exercise}
\label{ex:7.16}
(cont.) Show that $\varphi(x) = (0\, 1\, 2\, 3\, 4\, 5\, 6)$, that $\varphi(y) = (1\, 4\, 2)(3\, 5\, 6)$, and that $\varphi(z) = z_i$ for some $i$, where $z_1 = (\infty\, 0)(1\, 3)(2\, 5)(4\, 6)$, $z_2 = (\infty\, 0)(1\, 5)(2\, 6)(3\, 4)$, and $z_3 = (\infty\, 0)(1\, 6)(2\, 3)(4\, 5)$.
\end{exercise}
\begin{solution}
    The action is left multiplication on the eight left cosets of $H$. With the labelling above, $\infty$ denotes the coset $H$, while $i$ denotes the coset $x^izH$ for $0\leq i\leq 6$. Since $x\in H$, the element $x$ fixes $H$. Also, for $0\leq i\leq 6$,
    \[
        x\cdot x^izH=x^{i+1}zH,
    \]
    where the exponent is read modulo $7$. Hence
    \[
        \varphi(x)=(0\,1\,2\,3\,4\,5\,6).
    \]
    Similarly, $y\in H$, so $y$ fixes $H$. From Exercise~\ref{ex:7.13}, we have $yxy^{-1}=x^4$, and hence
    \[
        yx^i=x^{4i}y
    \]
    for every $i$. From Exercise~\ref{ex:7.14}, we also have $yzy=z$, so $yz=zy^{-1}$. Since $y^{-1}\in H$, this gives $yzH=zH$. Therefore
    \[
        y\cdot x^izH=yx^izH=x^{4i}yzH=x^{4i}zH.
    \]
    Thus $\varphi(y)$ sends $i$ to $4i$ modulo $7$, while fixing $\infty$ and $0$. Hence
    \[
        \varphi(y)=(1\,4\,2)(3\,5\,6).
    \]
    It remains to determine the possible form of $\varphi(z)$. Write $t=\varphi(z)$ and $\sigma=\varphi(y)$. Since $z^2=1$, we have $t^2=1$. Also,
    \[
        z\cdot H=zH
        \quad\text{and}\quad
        z\cdot zH=H,
    \]
    so $t$ contains the transposition $(\infty\,0)$. Since $\varphi(G)\leq A_8$, the permutation $t$ is even. Finally, Exercise~\ref{ex:7.14} gives $zyz^{-1}=y^{-1}$, so
    \[
        t\sigma t^{-1}=\sigma^{-1}.
    \]
    The two non-trivial cycles of $\sigma$ have supports $\{1,4,2\}$ and $\{3,5,6\}$. The last conjugation relation forces $t$ either to preserve these two supports or to interchange them. It cannot preserve both supports: on each support it would have to reverse a $3$-cycle, giving one transposition on each support, and together with $(\infty\,0)$ this would make $t$ odd. Hence $t$ interchanges the two supports.\\
    Thus $t(1)$ is one of $3,5,6$. Since $t\sigma=\sigma^{-1}t$, the image of $1$ determines the rest:
    \[
        \begin{array}{c|c|c|c}
            t(1) & t(4) & t(2) & t\\
            \hline
            3 & 6 & 5 & (\infty\,0)(1\,3)(2\,5)(4\,6)\\
            5 & 3 & 6 & (\infty\,0)(1\,5)(2\,6)(3\,4)\\
            6 & 5 & 3 & (\infty\,0)(1\,6)(2\,3)(4\,5).
        \end{array}
    \]
    Therefore $\varphi(z)=z_i$ for some $i\in\{1,2,3\}$.
\end{solution}

\begin{exercise}
\label{ex:7.17}
(cont.) Let $L_i = \langle \varphi(x), \varphi(y), z_i \rangle \leq A_8$ for $i = 1, 2, 3$. Show that $L_1 = L_2 = L_3$. Conclude that $G$ is uniquely determined up to isomorphism.
\end{exercise}
\begin{solution}
    Put
    \[
        a=\varphi(x)=(0\,1\,2\,3\,4\,5\,6)
        \quad\text{and}\quad
        b=\varphi(y)=(1\,4\,2)(3\,5\,6).
    \]
    Then $L_i=\langle a,b,z_i\rangle$. We first compare the three possible choices of $z_i$. Since $b$ fixes $\infty$ and $0$, direct conjugation gives
    \[
        bz_1b^{-1}
        =(\infty\,0)(4\,5)(1\,6)(2\,3)
        =z_3
    \]
    and
    \[
        b^{-1}z_1b
        =(\infty\,0)(2\,6)(4\,3)(1\,5)
        =z_2.
    \]
    Thus $z_2,z_3\in L_1$, and hence
    \[
        L_2=\langle a,b,z_2\rangle\leq L_1,
        \qquad
        L_3=\langle a,b,z_3\rangle\leq L_1.
    \]
    Conversely, the same two identities give
    \[
        z_1=bz_2b^{-1}
        \quad\text{and}\quad
        z_1=b^{-1}z_3b.
    \]
    Hence $z_1\in L_2$ and $z_1\in L_3$, so $L_1\leq L_2$ and $L_1\leq L_3$. Therefore
    \[
        L_1=L_2=L_3.
    \]
    It remains to connect this with the original group $G$. By Exercise~\ref{ex:7.15}, we have $G=\langle x,y,z\rangle$. Applying $\varphi$ gives
    \[
        \varphi(G)=\langle \varphi(x),\varphi(y),\varphi(z)\rangle.
    \]
    By Exercise~\ref{ex:7.16}, the element $\varphi(z)$ is one of $z_1,z_2,z_3$. Hence
    \[
        \varphi(G)=L_i
    \]
    for some $i\in\{1,2,3\}$. Since $L_1=L_2=L_3$, the image $\varphi(G)$ is the fixed subgroup $L_1\leq A_8$, independent of the possible choice of $\varphi(z)$. Finally, $\varphi$ is a monomorphism, so
    \[
        G\cong \varphi(G)=L_1.
    \]
    Thus $G$ is uniquely determined up to isomorphism.
\end{solution}

\end{document}
